\pdfvariable suppressoptionalinfo 512
\pdfvariable trailerid {[<C3562026090300000000000000000000><C3562026090300000000000000000000>]}
\documentclass[10pt]{article}
\usepackage[margin=0.65in]{geometry}
\usepackage{amsmath,amssymb,amsthm,booktabs,microtype,xurl}
\usepackage{unicode-math}
\setmainfont{Latin Modern Roman}
\setmathfont{Latin Modern Math}
\providecommand{\CRevisionRound}{2}
\newtheorem{theorem}{Theorem}
\newtheorem{corollary}{Corollary}
\newcommand{\Tr}{\operatorname{Tr}}
\newcommand{\T}{\mathbb T}
\title{An Exact Chern-Phase Atlas for the QWZ Thouless Pump}
\author{HCS-C356 / HEN-O340}
\date{3 September 2026}
\begin{document}
\maketitle
\begin{abstract}
For a two-band pump on the Bloch torus we prove the exact spectrum, direct gap,
all lower-band Chern phases, and every gap-closing Dirac jump.  The transported
charge statement is made only in the filled-band adiabatic limit.  Independent
symbolic, degree, and lattice-gauge receipts audit the convention-sensitive
signs.  This revision is the \textbf{exact Bloch-gap owner}.
\end{abstract}

\section{Model and complete gap atlas}
Orient $\T^2$ by $dk\wedge d\tau$ and define
\begin{equation}\label{eq:model}
 H_m(k,\tau)=d_m(k,\tau)\cdot\sigma,\qquad
 d_m=(\sin k,\sin\tau,m+\cos k+\cos\tau).
\end{equation}
Here $m\in\mathbb R$ and both angular variables have period $2\pi$.  Whenever
$d_m$ is nonzero, let $n=d_m/|d_m|$ and $P_-=(I-n\cdot\sigma)/2$.

\begin{theorem}[Spectrum, gap, and lower-band phase]\label{thm:main}
The fibre spectrum is $\{-|d_m|,+|d_m|\}$.  The family is gapped exactly for
$m\notin\{-2,0,2\}$, with direct gap
\[
 G(m)=2\min\{|m+2|,|m|,|m-2|\}.
\]
On a gapped chamber $P_-$ is a smooth rank-one complex line bundle.  With
\begin{equation}\label{eq:c1}
 c_1(P_-)=\frac1{2\pi i}\int_{\T^2}
 \Tr\!\left(P_-[\partial_kP_-,\partial_\tau P_-]\right)dk\,d\tau,
\end{equation}
its Chern number is
\[
c_1=\begin{cases}
0,&m<-2,\\ -1,&-2<m<0,\\ +1,&0<m<2,\\ 0,&m>2.
\end{cases}
\]
\end{theorem}

\begin{proof}[Proof of the spectral and gap assertions]
Pauli multiplication gives $H_m^2=|d_m|^2I$.  With $x=\cos k$ and
$y=\cos\tau$,
\begin{equation}\label{eq:norm}
 |d_m|^2=m^2+2+2m(x+y)+2xy.
\end{equation}
The right side is affine in each variable separately, so its minimum on
$[-1,1]^2$ occurs at a corner.  The four corner values are
$(m+2)^2,m^2,m^2,(m-2)^2$.  This proves both the closing locus and the stated
gap.  Smooth functional calculus then gives the displayed lower projector.
The Chern computation is proved below in revision one.
\end{proof}

\ifnum\CRevisionRound>0
\section{Projector sign, degree, and Dirac walls}
This revision is the \textbf{analytic Dirac-degree owner}.  Direct Pauli
algebra fixes the convention-sensitive sign:
\begin{equation}\label{eq:pauli}
 \Tr\!\left(P_-[\partial_kP_-,\partial_\tau P_-]\right)
 =-\frac{i}{2}n\cdot(\partial_kn\times\partial_\tau n).
\end{equation}
Consequently \eqref{eq:c1} is minus the Brouwer degree of $n:\T^2\to S^2$.
For completeness, a direct differentiation before normalization gives
\begin{equation}\label{eq:numerator}
d_m\cdot(\partial_kd_m\times\partial_\tau d_m)
=\cos k+\cos\tau+m\cos k\cos\tau.
\end{equation}

Choose the north pole as regular value away from the walls.  Its possible
preimages are the four points with $\sin k=\sin\tau=0$.  A point contributes
exactly when its mass $M=m+\cos k_0+\cos\tau_0$ is positive, and its local
degree is $\chi(k_0,\tau_0)=\cos k_0\cos\tau_0$.  Thus
$\deg(n)=\sum\chi\mathbf1_{\{M>0\}}$.  Substituting
$\mathbf1_{\{M>0\}}=(1+\operatorname{sgn}M)/2$ and using
$\sum\chi=0$ yields
\begin{equation}\label{eq:masssum}
c_1=-\frac12\bigl[\operatorname{sgn}(m+2)
-2\operatorname{sgn}(m)+\operatorname{sgn}(m-2)\bigr].
\end{equation}
This proves the remaining part of Theorem~\ref{thm:main} without numerical
integration.  As $m$ increases, the complete wall ledger is
\begin{center}
\begin{tabular}{cccc}
\toprule
mass & zero $(k_0,\tau_0)$ & $\chi$ & jump of $c_1$\\
\midrule
$-2$ & $(0,0)$ & $+1$ & $-1$\\
$0$ & $(\pi,0)$ & $-1$ & $+1$\\
$0$ & $(0,\pi)$ & $-1$ & $+1$\\
$2$ & $(\pi,\pi)$ & $+1$ & $-1$\\
\bottomrule
\end{tabular}
\end{center}
The simultaneous pair at $m=0$ is essential: its total jump is two.
\fi

\ifnum\CRevisionRound>1
\section{Adiabatic transport and exact boundaries}
This revision is the \textbf{adiabatic scope and route firewall}.
\begin{corollary}[Filled-band Thouless transport]
Let $\tau$ increase from $0$ to $2\pi$, assign unit positive particle charge,
and define positive current by $J=\partial_kH_m$.  Assume a gapped chamber, a
completely filled lower band, and an adiabatic traversal.  With the orientation
fixed in \eqref{eq:c1}, the transported positive-particle charge per cycle
equals $c_1(P_-)$.
\end{corollary}
Indeed, for
\[
\Omega_{k\tau}=i\Tr P_-[\partial_kP_-,\partial_\tau P_-],
\qquad
Q=-\frac1{2\pi}\int_{\T^2}\Omega_{k\tau}\,dk\,d\tau
  =\frac1{2\pi i}\int_{\T^2}\Tr P_-[\partial_kP_-,\partial_\tau P_-]
  =c_1(P_-).
\]
For electronic charge the result is multiplied by $-e$.
This is an adiabatic-limit statement.  At finite driving rate, interband
transitions can occur; no exact finite-driving-rate quantization is claimed.

At $m=-2,0,2$ the lower projector is undefined at the zeroes in the table and
no gapped-bundle Chern number is assigned.  Reversing the torus orientation,
or sending $\tau\mapsto-\tau$, reverses the integer.  Complex conjugation
satisfies
\[
K H_m(k,\tau)K=H_m(k,-\tau),
\]
because only $\sigma_y$ is imaginary.  For $m>2$ (respectively $m<-2$), the
$z$ component is everywhere positive (respectively negative), and the field
contracts without closing the gap to the north (respectively south) pole.

\section{Evidence, provenance, and claim boundary}
The exact evidence records the four corner values, chamber signs, and Dirac
ledger.  A separate discrete lattice-gauge lane on several grids reproduces
the four integers, but is only regression evidence; the analytic degree proof
owns the theorem.  Strict parsers, repaired-hash hostile mutations, isolated
byte replay, and deterministic PDF builds audit the release.

Thouless established quantized adiabatic particle transport~\cite{Thouless};
Qi, Wu, and Zhang give the two-band lattice-Dirac lineage~\cite{QWZ}.  We cite
lineage, not priority, and rederive every normalization and sign used here.
Nearby packages treat finite SSH chains, the Dirac-monopole sphere and its
magnetic spectrum, and kicked rotors; none owns this smooth two-parameter
Bloch-bundle atlas.

The Chern integer is a source-local natural quantization, but the Chern integer
is not a rational-prime carrier.  There is no primitive-orbit determinant,
target functional equation, target zero match, or Hilbert--P\'olya operator.
Thus Route A is rejected as
$(A0_{\rm FAIL},A1_{\rm FAIL},A2_{\rm FAIL},A3_{\rm FAIL},
A4_{\rm NATURAL\_QUANTIZATION})$, and Route B remains locked under
\texttt{NO\_BAD\_EULER\_OR\_ROOT\_NUMBER}.

\begin{thebibliography}{9}
\small
\bibitem{Thouless}
D.~J. Thouless, \emph{Quantization of particle transport}, Phys. Rev. B 27
(1983), 6083--6087. \url{https://doi.org/10.1103/PhysRevB.27.6083}.
\bibitem{QWZ}
X.-L. Qi, Y.-S. Wu, and S.-C. Zhang, \emph{Topological quantization of the
spin Hall effect in two-dimensional paramagnetic semiconductors}, Phys. Rev.
B 74 (2006), 085308. \url{https://doi.org/10.1103/PhysRevB.74.085308}.
\end{thebibliography}
\fi
\end{document}
